\Needspace{18\baselineskip}
\section{Exact equations on products with zero lower mass}
\label{sec:zero-lower-mass-product-equations}

Positive lower mass guarantees solvability for all measures, but an individual density can have a positive threshold and an exact equation even when the lower mass vanishes. We construct such examples on products $X_0\times Y$ of a Hermitian manifold with $v_-(h_0)=0$ and a closed K\"ahler factor. On $Y=\mathbb P^k$, a potential with a sublogarithmic pole, pulled back from the projective factor, solves an explicit equation whose density is integrable but belongs to no $L^p$ with $p>1$. The coefficient is computed from the normalization scalars of bounded truncations, the equation follows from the product structure of the Monge--Amp\`ere measure, and the product background again has zero lower mass.

Throughout this section, $(X_0,h_0)$ is a connected compact Hermitian manifold of dimension $m\geq1$, $(Y,\kappa)$ is a connected compact K\"ahler manifold of dimension $k\geq1$ normalized by $\int_Y\kappa^k=1$, and $X=X_0\times Y$, a compact complex manifold of dimension $N=m+k$ with the Hermitian background $H=h_0+\kappa$. Forms and functions on the factors are identified with their pullbacks to $X$. Put
\[
 V_0=\int_{X_0}h_0^m,\qquad C=\binom Nm,\qquad
 \mu=V_0^{-1}h_0^m\wedge\kappa^k,\qquad R=\frac{\mu}{\mathrm{d}V},
\]
where $\mathrm{d}V$ is any smooth positive probability volume form on $X$. Then $\mu$ is a smooth positive probability volume form, and $R$ is smooth and bounded above and below by positive constants. The normalization scalar $c_H$ and the threshold $b_H$ are those of \cref{sec:lower-mass-solvability,sec:individual-density-subsolutions}, now on the $N$-dimensional manifold $X$ with background $H$.

\Needspace{6\baselineskip}
\begin{lemma}\label{lem:split-product-normalization}\index{Normalization scalar!split products}
\begin{enumerate}
\item For every $v\in\PSH(Y,\kappa)$,
\begin{equation}\label{eq:split-product-ma}
 \mathrm{MA}^{\mathrm{ext}}_H(v)=Ch_0^m\wedge\mathrm{MA}^{\mathrm{ext}}_\kappa(v).
\end{equation}
\item For every nonzero nonnegative $g\in L^1(Y,\kappa^k)$, with $M_g=\int_Yg\kappa^k$,
\[
 b_H(Rg)=\frac{CV_0}{M_g},
\]
and $c_H(Rg)=CV_0/M_g$ if $g$ is bounded.
\item If $v_-(h_0)=0$, then $v_-(H)=0$ and $H$ is not closed.
\end{enumerate}
\end{lemma}
\begin{proof}
(1) Forms and currents of bidegree $(j,j)$ pulled back from $X_0$ vanish for $j>m$, and those pulled back from $Y$ vanish for $j>k$. For bounded $v$, \cref{prop:bounded-products,lem:bounded-products-submersion} and the binomial expansion therefore give $\mathrm{MA}_H(v)=Ch_0^m\wedge\mathrm{MA}_\kappa(v)$. For general $v$, apply this to $\max(v,-j)$, multiply by the pullback of $\mathbf1_{\{v>-j\}}$, and let $j\to\infty$ as in the definition \eqref{eq:extended-ma} of the extended measure.

(2) Suppose first that $g$ is bounded. \Cref{thm:kn-continuous} on $Y$, with reference volume $\kappa^k$, gives a continuous $v\in\PSH(Y,\kappa)$ and $c>0$ with $\mathrm{MA}_\kappa(v)=cg\kappa^k$. Since $\kappa$ is closed, Stokes' theorem gives $\int_Y\mathrm{MA}_\kappa(v)=1$, hence $c=1/M_g$, and (1) gives
\[
 \mathrm{MA}_H(v)=\frac{C}{M_g}h_0^m\wedge g\kappa^k
 =\frac{CV_0}{M_g}Rg\,\mathrm{d}V.
\]
Uniqueness of the scalar gives $c_H(Rg)=CV_0/M_g$. For general $g$, write $a=\min_XR$ and $A=\max_XR$. The truncations of $Rg$ that define $b_H(Rg)$ in \eqref{eq:individual-truncated-normalizer} satisfy
\[
 R\min(g,j/A)\leq\min(Rg,j)\leq R\min(g,j/a).
\]
By the bounded case and the monotonicity of the scalar in \cref{lem:bounded-normalizer-subsolution},
\[
 \frac{CV_0}{M_{\min(g,j/a)}}\leq c_H(\min(Rg,j))
                         \leq\frac{CV_0}{M_{\min(g,j/A)}}.
\]
By monotone convergence, both bounds tend to $CV_0/M_g$.

(3) For bounded $w\in\PSH(X_0,h_0)$, the same expansion gives $\mathrm{MA}_H(w)=C\,\mathrm{MA}_{h_0}(w)\wedge\kappa^k$, a measure of mass $C\int_{X_0}\mathrm{MA}_{h_0}(w)$. Pulling back bounded potentials whose masses on $X_0$ tend to zero gives $v_-(H)=0$. If $h_0$ were closed, Stokes' theorem would give every bounded $h_0$-plurisubharmonic function the mass $V_0>0$, contradicting $v_-(h_0)=0$. Hence $h_0$, and therefore $H$, is not closed.
\end{proof}

Until \cref{subsec:closed-factor-equations}, let $Y=\mathbb P^k$, with $\kappa$ the Fubini--Study form normalized by $\int_Y\kappa^k=1$. For $0<\alpha<1$ define, on $Y$,
\[
 \psi([z_0:\cdots:z_k])=\log\frac{\sum_{i=1}^k|z_i|^2}{\sum_{i=0}^k|z_i|^2},
 \qquad u_\alpha=1-(1-\psi)^\alpha.
\]
The function $\psi$ is smooth on $Y\setminus\{o\}$, where $o=[1:0:\cdots:0]$, and tends to $-\infty$ at $o$; moreover $\psi\leq0$, with equality exactly on the hyperplane $\{z_0=0\}$. Thus $u_\alpha\leq0$, and the pole set of $u_\alpha$ is $\{o\}$.

\Needspace{6\baselineskip}
\begin{lemma}\label{lem:projective-sublogarithmic-density}\index{Monge--Ampere density@Monge--Amp\`ere density!sublogarithmic pole}
The function $u_\alpha$ is $\kappa$-plurisubharmonic with supremum zero, and
\[
 \mathrm{MA}^{\mathrm{ext}}_\kappa(u_\alpha)=q_\alpha\kappa^k,
 \qquad \int_Yq_\alpha\kappa^k=1.
\]
The density $q_\alpha$ is smooth and strictly positive on $Y\setminus\{o\}$, belongs to $L^1(\kappa^k)$, and belongs to no $L^p(\kappa^k)$ with $p>1$. If we set $q_\alpha(o)=+\infty$, then every truncation $\min(q_\alpha,s)$, $s>0$, is strictly positive and continuous on $Y$.
\end{lemma}
\begin{proof}
In homogeneous coordinates, $\kappa+\ddc\psi=\ddc\log\sum_{i=1}^k|z_i|^2\geq0$ on $Y\setminus\{o\}$. The function $\chi(t)=1-(1-t)^\alpha$, $t\leq0$, is increasing and convex, with $0<\chi'\leq\alpha$. Hence, off $o$,
\[
 \kappa+\ddc u_\alpha\geq(1-\chi'(\psi))\kappa+\chi'(\psi)(\kappa+\ddc\psi)
 \geq(1-\alpha)\kappa.
\]
Since $u_\alpha$ is bounded above, it extends across $o$ as a $\kappa$-plurisubharmonic function, with value $-\infty$ at $o$. On the hyperplane $\{z_0=0\}$ we have $\psi=0$, hence $u_\alpha=0$, so $\sup_Yu_\alpha=0$.

To compute the measure, work in the affine chart $z_0=1$ and put $r=\sum_{i=1}^k|z_i|^2$. On this chart $\kappa=\ddc\log(1+r)$ and $\kappa+\ddc u_\alpha=\ddc F$, where $s(r)=1-\log(r/(1+r))$ and
\[
 F(r)=\log(1+r)+1-s(r)^\alpha,\qquad
 Q(r)=rF'(r)=\frac{r+\alpha s(r)^{\alpha-1}}{1+r}.
\]
For a radial function, the Levi form has the tangential eigenvalue $F'(r)=Q(r)/r$, with multiplicity $k-1$, and the radial eigenvalue
\[
 Q'(r)=\frac{1-\alpha s(r)^{\alpha-1}+\alpha(1-\alpha)r^{-1}s(r)^{\alpha-2}}{(1+r)^2}>0.
\]
Integrating in the radial variable gives $\int_{\{a<r<b\}}(\ddc F)^k=Q(b)^k-Q(a)^k$; the constant of proportionality is fixed by the case $F=\log(1+r)$, for which $(\ddc F)^k=\kappa^k$ has mass one on the chart and $Q(r)=r/(1+r)$. Since again $Q(0)=0$ and $Q(\infty)=1$, the punctured chart has mass one. The complement of the chart is a hyperplane, which has zero volume, and the non-pluripolar measure does not charge $o$. Hence $\int_Yq_\alpha\kappa^k=1$.

Relative to $\kappa^k$, whose tangential and radial eigenvalues are $1/(1+r)$ and $1/(1+r)^2$, the density on this chart is
\[
 q_\alpha(r)=(1+r)^{k+1}(Q(r)/r)^{k-1}Q'(r)
 \asymp r^{-k}(1-\log r)^{-1-k(1-\alpha)}\quad(r\downarrow0).
\]
Near $r=0$ the smooth volume is comparable to $r^{k-1}\mathrm{d}r$ in the radial variable, so $q_\alpha$ is integrable. For $p>1$, the substitution $L=1-\log r$ turns the integral of $q_\alpha^p$ into the integral of $\mathrm{e}^{k(p-1)L}$ times a power of $L$, which diverges. The same asymptotics show that $q_\alpha(r)\to+\infty$ uniformly as $r\downarrow0$, so every truncation $\min(q_\alpha,s)$ is continuous at $o$. Finally, $q_\alpha\geq(1-\alpha)^k>0$ by the lower bound $\kappa+\ddc u_\alpha\geq(1-\alpha)\kappa$.
\end{proof}

\Needspace{6\baselineskip}
\begin{theorem}\label{thm:zero-lower-mass-product-equation}\index{Monge--Ampere equation@Monge--Amp\`ere equation!zero lower mass products}
For every $0<\alpha<1$, the function $f=Rq_\alpha$ is a probability density in $L^1(\mathrm{d}V)$, belonging to no $L^p(\mathrm{d}V)$ with $p>1$, and
\[
 b_H(f)=CV_0>0,\qquad
 \mathrm{MA}^{\mathrm{ext}}_H(u_\alpha)=b_H(f)f\,\mathrm{d}V.
\]
The potential $u_\alpha$, pulled back to $X$, has supremum zero and finite mass $CV_0$. It belongs to $\mathcal E_{\mathrm{env}}(X,H,0)$, and $P_{\delta H}(u_\alpha)\not\equiv-\infty$ for every $0<\delta<1$. If $v_-(h_0)=0$, then $H$ is not closed and $v_-(H)=0$.
\end{theorem}
\begin{proof}
\emph{Step 1: the equation.} By \cref{lem:projective-sublogarithmic-density}, $\int_Xf\,\mathrm{d}V=\int_Xq_\alpha\mu=\int_Yq_\alpha\kappa^k=1$. Part (2) of \cref{lem:split-product-normalization}, with $g=q_\alpha$ and $M_g=1$, gives $b_H(f)=CV_0$, and part (1) gives
\[
 \mathrm{MA}^{\mathrm{ext}}_H(u_\alpha)=Ch_0^m\wedge q_\alpha\kappa^k
 =CV_0\,q_\alpha\mu=CV_0f\,\mathrm{d}V,
\]
a measure of mass $CV_0$. Since $R$ is bounded above and below by positive constants, $f$ has the integrability properties of $q_\alpha$. The last assertion is part (3) of the lemma.

\emph{Step 2: the envelopes.} Choose $\alpha<\beta<1$. For $b\geq1$ and $\varepsilon>0$ with $\varepsilon\beta\leq1$, the number
\[
 K_b=\sup_{s\geq1}\{b s^\alpha-\varepsilon s^\beta\}-b
\]
is finite, since $\beta>\alpha$. The function $g_b=-\varepsilon(1-\psi)^\beta-K_b$ on $Y$ satisfies $\kappa+\ddc g_b\geq(1-\varepsilon\beta)\kappa\geq0$, by the chain-rule computation in the proof of \cref{lem:projective-sublogarithmic-density}, applied to the increasing convex function $t\mapsto-\varepsilon(1-t)^\beta$, whose derivative is at most $\varepsilon\beta$ for $t\leq0$. With $L=1-\psi\geq1$, the inequality $g_b\leq bu_\alpha$ reads $bL^\alpha-\varepsilon L^\beta-b\leq K_b$, so it holds on $Y\setminus\{o\}$, and also at $o$, where both sides are $-\infty$. Thus the pullback of $g_b$ belongs to $\PSH(X,H)$ and lies below $bu_\alpha$. Since $b\geq1$ is arbitrary, $u_\alpha\in\mathcal E_{\mathrm{env}}(X,H,0)$ by \cref{def:relative-envelope-class}. For $0<\delta<1$, take $b=1$ and $\varepsilon\beta\leq\delta$. Then $g_1$ is $\delta H$-plurisubharmonic and lies below $u_\alpha$, so $P_{\delta H}(u_\alpha)\geq g_1$ is not identically $-\infty$.
\end{proof}

\Needspace{6\baselineskip}
\begin{corollary}\label{cor:elliptic-fivefold-exact-threshold}\index{Elliptic fibration!exact product equation}
Let $X_0$ be the elliptic threefold of \cref{thm:vanishing-lower-mass}, and let $h_0$ be the balanced metric of \cref{prop:lower-mass-balanced-metric}. On $X_0\times\mathbb P^2$, with background $H=h_0+\kappa$, the densities $f=Rq_\alpha$ of the preceding theorem satisfy
\[
 v_-(H)=0,\qquad b_H(f)=240\pi,
 \qquad\mathrm{MA}^{\mathrm{ext}}_H(u_\alpha)=240\pi f\,\mathrm{d}V.
\]
\end{corollary}
\begin{proof}
\Cref{thm:vanishing-lower-mass} gives $v_-(h_0)=0$, and \cref{prop:lower-mass-balanced-metric} gives $V_0=24\pi$. The product theorem applies with $m=3$, $k=2$ and $C=\binom53=10$.
\end{proof}

On these products, the radial flux also identifies the limit of the supremum-normalized exponential solutions as the exponent tends to zero (\cref{thm:radial-integrable-selection}); a first moment of the limiting measure determines whether the solutions shifted by $-\lambda^{-1}\log b_H(f)$ have a finite limit (\cref{prop:radial-scalar-centering}). \Cref{qu:individual-equation-threshold} asks for an exact equation for every integrable density with positive threshold.

\subsection{Closed factors and arbitrary integrable densities}
\label{subsec:closed-factor-equations}

The preceding examples use an explicit potential on projective space. On a closed K\"ahler factor, Stokes' theorem instead bounds the non-pluripolar mass, and this bound makes every threshold subsolution on the factor exact. This gives equations for arbitrary integrable densities on that factor, with no assumption that they are bounded below.

In this subsection $(Y,\kappa)$ is again an arbitrary connected compact K\"ahler manifold with $\int_Y\kappa^k=1$, with the notation $X$, $N$, $H$, $C$, $V_0$, $\mu$ and $R$ introduced at the beginning of the section, and we put $c_0=CV_0$.

\Needspace{12\baselineskip}
\begin{theorem}\label{thm:closed-factor-integrable-equation}\index{Monge--Ampere equation@Monge--Amp\`ere equation!closed factor}\index{Normalization scalar!closed factor}
For every nonzero nonnegative $g\in L^1(Y,\kappa^k)$, with $M_g=\int_Yg\kappa^k$, one has
\[
 b_H(Rg)=\frac{c_0}{M_g}.
\]
There is $v\in\PSH(Y,\kappa)$ with $\sup_Yv=0$ and $\int_Y\mathrm{MA}^{\mathrm{ext}}_\kappa(v)=1$ whose pullback $u$ satisfies
\[
 \mathrm{MA}^{\mathrm{ext}}_H(u)
       =\frac{c_0}{M_g}Rg\,\mathrm dV,\qquad
 \int_X\mathrm{MA}^{\mathrm{ext}}_H(u)=c_0.
\]
\end{theorem}
\begin{proof}
The formula for $b_H(Rg)$ is part (2) of \cref{lem:split-product-normalization}. By \cref{prop:individual-normalizer-lp}, applied on $Y$ with the K\"ahler form $\kappa$ and the reference volume $\kappa^k$, $b_\kappa(g)=1/M_g$. Thus \cref{thm:individual-density-subsolution} on $Y$, followed by subtraction of a constant, gives $v\in\PSH(Y,\kappa)$ with supremum zero and
\[
 \mathrm{MA}^{\mathrm{ext}}_\kappa(v)\geq(g/M_g)\kappa^k.
\]
For $v_j=\max(v,-j)$, Stokes' theorem for the closed form $\kappa$ gives $\int_Y\mathrm{MA}_\kappa(v_j)=1$. Since $\mathrm{MA}^{\mathrm{ext}}_\kappa(v)$ is the increasing limit of the measures $\mathbf1_{\{v>-j\}}\mathrm{MA}_\kappa(v_j)$ in \eqref{eq:extended-ma}, it follows that $\int_Y\mathrm{MA}^{\mathrm{ext}}_\kappa(v)\leq1$. The prescribed lower measure has mass one, so equality holds both in mass and in the measure inequality. Part (1) of \cref{lem:split-product-normalization} now gives
\[
 \mathrm{MA}^{\mathrm{ext}}_H(u)
   =C h_0^m\wedge\mathrm{MA}^{\mathrm{ext}}_\kappa(v)
   =\frac{c_0}{M_g}Rg\,\mathrm dV,
\]
a measure of mass $c_0$.
\end{proof}

When $v_-(h_0)=0$, the background $H$ has zero lower mass and is not closed, by part (3) of \cref{lem:split-product-normalization}. For $Y=\mathbb P^k$ and $g=q_\alpha$, the potential $u_\alpha$ of \cref{thm:zero-lower-mass-product-equation} is such a solution. For the elliptic fivefold of \cref{cor:elliptic-fivefold-exact-threshold}, $c_0=240\pi$; thus, for every probability density $g$ on $\mathbb P^2$, the density $Rg$ has threshold $240\pi$ and admits a threshold solution of finite mass. The theorem concerns factor densities; it does not give an exact equation for an arbitrary density on $X$ or identify limits of normalized exponential solutions.

